\makeatletter
\@ifundefined{safeimage}{%
  \newcommand{\safeimage}[2]{%
    \def\safe@found{}%
    \IfFileExists{images/#1}{\def\safe@found{images/#1}}{%
      \IfFileExists{Images/#1}{\def\safe@found{Images/#1}}{%
        \IfFileExists{../images/#1}{\def\safe@found{../images/#1}}{%
          \IfFileExists{../Images/#1}{\def\safe@found{../Images/#1}}{%
            \IfFileExists{FSI_template_Memoire_final/images/#1}{\def\safe@found{FSI_template_Memoire_final/images/#1}}{%
              \IfFileExists{FSI_template_Memoire_final/Images/#1}{\def\safe@found{FSI_template_Memoire_final/Images/#1}}{}}}}}}%
    \ifx\safe@found\@empty
      \fbox{\parbox[c][4.5cm][c]{#2}{\centering Image manquante : \detokenize{#1}}}%
    \else
      \includegraphics[width=#2]{\safe@found}%
    \fi
  }%
}{}
\makeatother

\chapter{Résultats et Discussion}

\section{Introduction}

Ce chapitre présente et discute les résultats obtenus lors de la conception et de
l'évaluation de la plateforme \textsc{SAFE CONGO}, un système de surveillance
épidémiologique et d'alerte précoce destiné à la RDC. L'objectif est double :
d'une part, caractériser le comportement des séries épidémiologiques hebdomadaires
à travers une analyse exploratoire rigoureuse ; d'autre part, évaluer la capacité
des modèles d'apprentissage automatique à projeter le nombre de cas par maladie et
à alimenter un mécanisme d'alerte conforme aux standards de l'OMS et de la
SIMR/IDSR.

Les résultats rapportés ici sont issus directement des artefacts produits par le
pipeline de traitement et d'entraînement : le rapport d'analyse exploratoire, les
tableaux de performance par maladie, les matrices de confusion de sévérité et les
classements d'importance des variables.

\section{Protocole expérimental}

Le jeu de données nettoyé couvre la période allant du 14 janvier 2021 au
2 décembre 2023 et regroupe 2655 observations hebdomadaires réparties sur
25 maladies à déclaration obligatoire. Aucune date invalide ni année aberrante
n'a été conservée après nettoyage, garantissant l'intégrité temporelle des séries.

Pour chaque maladie, une série temporelle hebdomadaire du nombre total de cas est
construite, puis enrichie par un ensemble de variables explicatives à caractère
temporel :

\begin{itemize}
  \item les retards temporels d'ordre 1 à 4, qui capturent la mémoire récente de la série ;
  \item les moyennes mobiles sur 2, 3 et 4 semaines, qui réduisent le bruit hebdomadaire ;
  \item le rang temporel de la semaine, le mois et le trimestre, qui représentent la saisonnalité ;
  \item le taux de croissance, la volatilité sur 4 semaines et la tendance locale ;
  \item un code numérique de maladie, permettant un apprentissage différencié.
\end{itemize}

Une transformation logarithmique est appliquée à la cible afin d'atténuer la forte
asymétrie des distributions de cas. Six familles d'algorithmes de régression sont
comparées pour chaque maladie : la régression linéaire, la régression Ridge, les
forêts aléatoires, le Gradient Boosting, les k plus proches voisins et les machines
à vecteurs de support. La qualité prédictive est mesurée par le coefficient de
détermination $R^2$, l'erreur absolue moyenne, la racine de l'erreur quadratique
moyenne et l'erreur absolue moyenne en pourcentage, complétés par une validation
croisée.

Enfin, la sévérité épidémiologique est discrétisée en trois niveaux
(\textit{Faible}, \textit{Modéré}, \textit{Élevé}) à partir de seuils spécifiques
à chaque maladie, afin d'évaluer non seulement la précision numérique mais aussi la
pertinence opérationnelle des projections.

\section{Analyse exploratoire des données}

\subsection{Distribution et hétérogénéité}

L'analyse exploratoire confirme la nature fortement déséquilibrée des données
épidémiologiques. Sur l'ensemble du jeu de données nettoyé, la moyenne hebdomadaire
s'établit à 261 cas avec un écart-type de 548, illustrant une dispersion
considérable. La valeur médiane demeure modeste avec 64 cas, tandis que les
quantiles extrêmes atteignent 24\,919 cas au 90\up{e} centile, 73\,731 au
95\up{e} centile et 292\,226 au 99\up{e} centile. Cette asymétrie traduit la
coexistence de maladies endémiques à très forte charge épidémiologique et
d'événements rares mais intenses.

\begin{figure}[htbp]
  \centering
  \safeimage{eda_distribution_totalcas.png}{0.85\textwidth}
  \caption{Distribution du nombre hebdomadaire de cas après nettoyage.}
  \label{fig:eda-distribution}
\end{figure}

La figure~\ref{fig:eda-distribution} illustre cette asymétrie prononcée, justifiant
le recours à la transformation logarithmique appliquée lors de la phase de
modélisation.

\subsection{Dynamique temporelle}

Le classement des maladies par volume d'observations met en évidence la
prédominance des pathologies respiratoires et fébriles. L'infection respiratoire
aiguë est la maladie la mieux représentée avec 171 semaines d'observation, suivie
du paludisme confirmé avec 169 semaines, de la grippe avec 168 semaines, de la
fièvre typhoïde avec 166 semaines, de la rougeole avec 164 semaines et du paludisme
suspect avec 163 semaines. La figure~\ref{fig:eda-tendance} montre l'évolution
hebdomadaire du nombre total de cas, révélant des dynamiques saisonnières marquées
et des pics épidémiques récurrents.

\begin{figure}[htbp]
  \centering
  \safeimage{eda_tendance_hebdomadaire.png}{0.85\textwidth}
  \caption{Évolution hebdomadaire du nombre total de cas, toutes maladies confondues.}
  \label{fig:eda-tendance}
\end{figure}

\begin{figure}[htbp]
  \centering
  \safeimage{eda_top5_maladies.png}{0.85\textwidth}
  \caption{Évolution temporelle des cinq maladies les plus fréquentes entre 2021 et 2023.}
  \label{fig:eda-top5}
\end{figure}

\subsection{Structure des corrélations}

L'étude des corrélations entre variables numériques apporte un éclairage utile pour
la modélisation. Le nombre total de cas est fortement corrélé aux tranches d'âge
infantiles et scolaires : enfants de 5 à 15 ans ($r = 0{,}80$), enfants de
12 à 59 mois ($r = 0{,}77$), personnes de plus de 15 ans ($r = 0{,}75$). À
l'inverse, la corrélation avec le nombre de décès ($r = -0{,}01$) et la population
totale ($r = 0{,}11$) demeure très faible, confirmant que la charge de morbidité
n'est pas un simple reflet de la taille de la population.

\begin{figure}[htbp]
  \centering
  \safeimage{eda_correlation_heatmap.png}{0.78\textwidth}
  \caption{Matrice de corrélation des variables numériques du jeu de données nettoyé.}
  \label{fig:eda-corr}
\end{figure}

\section{Performance globale des modèles prédictifs}

Le tableau~\ref{tab:perf-globale} récapitule, pour chacune des 24 maladies
modélisées, le meilleur algorithme retenu et les métriques de performance associées,
classées par coefficient de détermination décroissant.

\begin{table}[htbp]
  \centering
  \resizebox{\textwidth}{!}{%
  \begin{tabular}{|l|l|c|c|c|c|c|}
    \hline
    \textbf{Maladie} & \textbf{Meilleur modèle} & $\boldsymbol{R^2}$ & \textbf{MAE} & \textbf{RMSE} & \textbf{MAPE} & \textbf{CV $R^2$} \\
    \hline
    Infection respiratoire aiguë & Gradient Boosting   & 0,986 & 2587  & 4488  & 17,5 & 0,909 \\
    Diarrhée sanglante           & Gradient Boosting   & 0,980 & 30    & 43    & 19,1 & 0,900 \\
    Paludisme confirmé           & Gradient Boosting   & 0,975 & 9038  & 16831 & 32,0 & 0,708 \\
    Paludisme suspect            & Gradient Boosting   & 0,967 & 15896 & 27731 & 13,1 & 0,876 \\
    PFA                          & Gradient Boosting   & 0,957 & 3     & 4     & 13,0 & 0,893 \\
    Grippe                       & Gradient Boosting   & 0,953 & 2557  & 4149  & 29,9 & 0,905 \\
    Fièvre typhoïde              & Gradient Boosting   & 0,953 & 2005  & 3622  & 48,2 & 0,884 \\
    Pneumonie                    & Random Forest       & 0,943 & 890   & 1399  & 36,2 & 0,814 \\
    Diarrhée aqueuse             & Gradient Boosting   & 0,939 & 577   & 1115  & 48,7 & 0,873 \\
    Rougeole                     & Gradient Boosting   & 0,934 & 240   & 359   & 30,4 & 0,895 \\
    COVID-19                     & Gradient Boosting   & 0,906 & 58    & 83    & 21,3 & 0,765 \\
    Fièvre jaune                 & SVR                 & 0,882 & 1     & 2     & 13,9 & 0,900 \\
    Tétanos maternel             & Régression linéaire & 0,865 & 0     & 0     & 21,2 & 0,656 \\
    MAPI légères                 & Random Forest       & 0,823 & 22    & 31    & 38,6 & 0,659 \\
    Monkeypox                    & SVR                 & 0,808 & 16    & 25    & 30,0 & 0,732 \\
    Méningite                    & Gradient Boosting   & 0,800 & 14    & 29    & 26,0 & 0,908 \\
    Choléra                      & Gradient Boosting   & 0,736 & 97    & 151   & 28,6 & 0,813 \\
    FHA                          & Ridge Regression    & 0,720 & 0,14  & 0,20  & 9,0  & 0,570 \\
    TNN                          & Gradient Boosting   & 0,720 & 3     & 4     & 23,9 & 0,744 \\
    Coqueluche                   & SVR                 & 0,672 & 4     & 5     & 40,6 & 0,562 \\
    Décès maternels              & Gradient Boosting   & 0,439 & 17    & 29    & 40,7 & 0,857 \\
    MAPI graves                  & Régression linéaire & 0,007 & 1     & 1     & 29,0 & 0,119 \\
    Rage                         & Ridge Regression    & 0,000 & 5     & 9     & 43,7 & 0,489 \\
    Peste                        & Régression linéaire & 0,000 & 1     & 1     & 36,5 & 0,069 \\
    \hline
  \end{tabular}%
  }
  \caption{Performance du meilleur modèle par maladie, triée par $R^2$ décroissant.}
  \label{tab:perf-globale}
\end{table}

Les résultats sont globalement très satisfaisants pour la majorité des maladies à
forte charge épidémiologique. Onze modèles atteignent un $R^2$ supérieur ou égal à
$0{,}90$, l'infection respiratoire aiguë culminant à $0{,}986$. Ces performances
indiquent que la dynamique hebdomadaire des principales maladies endémiques est
largement explicable à partir de l'historique récent et des composantes
saisonnières, validant l'approche par séries temporelles enrichies.

\section{Comparaison des familles d'algorithmes}

L'analyse comparative des six familles d'algorithmes fait ressortir une nette
domination du Gradient Boosting, retenu comme meilleur modèle pour 14 des
24 maladies. Le tableau~\ref{tab:familles} présente la répartition des modèles
vainqueurs.

\begin{table}[htbp]
  \centering
  \resizebox{0.55\textwidth}{!}{%
  \begin{tabular}{|l|c|}
    \hline
    \textbf{Famille d'algorithme} & \textbf{Nombre de maladies} \\
    \hline
    Gradient Boosting    & 14 \\
    SVR                  & 3  \\
    Régression linéaire  & 3  \\
    Random Forest        & 2  \\
    Ridge Regression     & 2  \\
    KNN                  & 0  \\
    \hline
  \end{tabular}%
  }
  \caption{Répartition du meilleur algorithme par maladie.}
  \label{tab:familles}
\end{table}

Cette prédominance s'explique par la capacité du Gradient Boosting à modéliser des
relations non linéaires et des interactions complexes entre les variables de retard
et de saisonnalité, tout en restant robuste face au bruit. Les forêts aléatoires
offrent une performance comparable sur certaines maladies à forte volatilité, tandis
que les modèles linéaires et SVR ne s'imposent que sur des maladies à très faibles
effectifs, où la simplicité du modèle limite le sur-apprentissage. Les k plus
proches voisins n'ont jamais été retenus comme meilleur modèle, leur sensibilité à
la dimensionnalité et à l'échelle des données les pénalisant systématiquement.

\section{Classification du niveau d'alerte}

\subsection{Matrices de confusion de sévérité}

Au-delà de la précision numérique, la finalité opérationnelle de la plateforme
impose d'évaluer la qualité de la classification du niveau de sévérité. Chaque
projection est convertie en un niveau \textit{Faible}, \textit{Modéré} ou
\textit{Élevé} à l'aide de seuils spécifiques à chaque maladie.

\begin{figure}[htbp]
  \centering
  \safeimage{matrices_confusion.png}{0.97\textwidth}
  \caption{Matrices de confusion de sévérité pour les principales maladies modélisées.}
  \label{fig:confusion}
\end{figure}

\subsection{Analyse détaillée du choléra}

Le tableau~\ref{tab:confusion-cholera} détaille la matrice de confusion obtenue pour
le choléra. Sur 19 semaines évaluées, le modèle classe correctement 13 observations,
soit une exactitude de 68,4\,\%. Les erreurs résiduelles se concentrent sur les
niveaux adjacents, sans aucune confusion directe entre les classes extrêmes
\textit{Faible} et \textit{Élevé}. Ce résultat est important en contexte d'alerte
précoce, car aucune flambée sévère n'est classée à tort comme bénigne.

\begin{table}[htbp]
  \centering
  \resizebox{0.70\textwidth}{!}{%
  \begin{tabular}{|l|c|c|c|}
    \hline
     & \textbf{Prédit Faible} & \textbf{Prédit Modéré} & \textbf{Prédit Élevé} \\
    \hline
    Réel Faible & 5 & 1 & 0 \\
    Réel Modéré & 2 & 3 & 0 \\
    Réel Élevé  & 0 & 3 & 5 \\
    \hline
  \end{tabular}%
  }
  \caption{Matrice de confusion du niveau de sévérité pour le choléra.}
  \label{tab:confusion-cholera}
\end{table}

\section{Importance des variables explicatives}

L'examen des importances de variables confirme le rôle central de la mémoire
temporelle des séries. Pour la plupart des maladies, les retards récents, les
moyennes mobiles et le rang temporel dominent le classement, devant les composantes
de tendance et de volatilité à quatre semaines. Cette structure est cohérente avec
la nature autorégressive des phénomènes épidémiques : le meilleur prédicteur du
nombre de cas d'une semaine donnée reste le niveau observé les semaines précédentes,
modulé par la dynamique saisonnière.

À titre d'illustration, pour la coqueluche le retard d'ordre 4 arrive en tête avec
une importance de 0,19, suivi du rang temporel avec 0,15 et de la moyenne mobile à
quatre semaines avec 0,10. Pour le COVID-19, le rang temporel domine avec 0,16,
devant la moyenne mobile à deux semaines avec 0,14 et la tendance locale avec 0,11.
La pertinence de ces variables justifie le choix d'un cadre de modélisation par
séries temporelles enrichies plutôt qu'une régression statique.

\section{Évaluation sur données non vues}

Une évaluation complémentaire sur un horizon de deux semaines a été conduite pour
l'ensemble des 23 maladies disposant de données de test. Les résultats sont
unanimes : tous les modèles obtiennent un verdict défavorable sur cet horizon court,
avec des $R^2$ négatifs sur les données hors-échantillon. Ce constat s'explique par
deux facteurs conjugués. D'une part, la brièveté de certaines séries limite la
généralisation à des configurations non observées lors de l'entraînement. D'autre
part, la prédiction à horizon fixe de deux semaines amplifie l'accumulation d'erreur
propre aux modèles autorégressifs. Ce résultat ne remet pas en cause la valeur des
modèles en production, qui opèrent en mode itératif en intégrant chaque nouvelle
observation dès qu'elle est saisie ; il délimite néanmoins honnêtement le domaine
de fiabilité des projections pures à moyen terme.

\section{Maladies à faible performance}

Quelques maladies présentent des performances limitées, voire nulles. La peste et la
rage affichent un $R^2$ de $0{,}000$, les MAPI graves un $R^2$ de $0{,}007$ et les
décès maternels un $R^2$ de $0{,}439$. Cette contre-performance s'explique
essentiellement par la rareté des événements et la brièveté des séries
correspondantes : ces pathologies ne comptent souvent que quelques dizaines de
semaines d'observation et présentent une forte proportion de valeurs nulles ou
unitaires. Dans ce régime de données clairsemées, aucun algorithme ne dispose d'un
signal suffisant pour apprendre une dynamique fiable. Ce constat délimite
honnêtement le domaine d'application de l'approche : la projection statistique n'est
pertinente que lorsque la profondeur d'historique est suffisante, critère formalisé
dans la plateforme par un seuil minimal de quatre semaines consécutives
d'observation.

\section{De la prédiction à l'alerte opérationnelle}

L'intégration des modèles au sein de la plateforme \textsc{SAFE CONGO} transforme
les projections en un dispositif d'alerte exploitable par les autorités sanitaires.
La figure~\ref{fig:app-saisie} présente l'espace administrateur, où l'agent encode
les cas hebdomadaires et obtient instantanément la projection du modèle ainsi que le
niveau d'alerte associé.

\begin{figure}[htbp]
  \centering
  \safeimage{app_saisie_admin.png}{0.92\textwidth}
  \caption{Interface de saisie et de projection dans l'espace administrateur.}
  \label{fig:app-saisie}
\end{figure}

Le niveau de sévérité est affiné en quatre paliers conformes à la classification
OMS/SIMR : \textsc{Faible}, \textsc{Modérée}, \textsc{Haute} et \textsc{Critique}.
Cette classification combine un critère sur le nombre de cas et un critère sur le
taux de croissance hebdomadaire, avec des seuils de 25\,\%, 50\,\% et 100\,\%.
Certaines maladies à tolérance zéro, telles que les fièvres hémorragiques aiguës,
déclenchent une alerte critique dès le premier cas déclaré, conformément au
Règlement Sanitaire International. Côté autorité sanitaire, ces alertes sont
centralisées dans un tableau de bord dédié, présenté à la figure~\ref{fig:app-alertes}.

\begin{figure}[htbp]
  \centering
  \safeimage{app_alertes_autorite.png}{0.92\textwidth}
  \caption{Tableau de bord des alertes côté autorité sanitaire.}
  \label{fig:app-alertes}
\end{figure}

Plusieurs garde-fous renforcent la fiabilité du dispositif. Les projections sont
plafonnées à dix fois le maximum historique récent, interdisant toute extrapolation
irréaliste. Lorsqu'aucune projection n'est disponible, faute d'un historique d'au
moins quatre semaines consécutives ou de modèle entraîné pour la zone considérée,
l'alerte repose exclusivement sur les cas observés sur le terrain et s'accompagne
d'une justification explicite transmise à l'autorité, précisant la raison de
l'absence de prévision.

\begin{figure}[htbp]
  \centering
  \safeimage{app_justification.png}{0.85\textwidth}
  \caption{Exemple d'alerte sans projection IA accompagnée de sa justification automatique.}
  \label{fig:app-justification}
\end{figure}

Cette transparence évite toute interprétation erronée et garantit la traçabilité de
chaque alerte émise.

\section{Discussion}

Les résultats obtenus démontrent la faisabilité d'un système de surveillance
prédictif opérationnel pour la RDC. La principale force de l'approche réside dans
la combinaison d'un pipeline reproductible, d'une sélection automatique du meilleur
modèle par maladie et d'une logique d'alerte alignée sur les standards
internationaux. Les performances élevées sur les maladies à forte charge
épidémiologique, avec onze modèles dépassant $R^2 = 0{,}90$, laissent entrevoir un
réel potentiel d'aide à la décision pour l'anticipation des flambées.

Plusieurs limites doivent néanmoins être soulignées. La période d'apprentissage,
restreinte aux années 2021 à 2023, demeure relativement courte au regard de la
variabilité interannuelle des épidémies. Les maladies rares restent difficilement
modélisables avec les approches classiques et nécessiteraient des méthodes
spécifiques aux données clairsemées. La granularité spatiale actuelle, centrée sur
la zone de santé, gagnerait à être affinée pour mieux saisir l'hétérogénéité
intraprovinciale. Enfin, l'absence de variables exogènes, notamment les données
météorologiques, la mobilité humaine et la couverture vaccinale, constitue une piste
d'amélioration prioritaire.

Ces limites tracent des perspectives claires : enrichissement continu des données,
intégration de covariables exogènes, déploiement de modèles adaptés aux faibles
effectifs et extension du dispositif à d'autres pays de la sous-région.

\section{Synthèse}

En définitive, ce chapitre a montré que la majorité des maladies à déclaration
obligatoire peuvent être projetées avec une précision élevée à partir de leur
historique hebdomadaire, le Gradient Boosting s'imposant comme l'algorithme de
référence pour 14 maladies sur 24. La transformation logarithmique de la cible, la
construction de variables de retard et de saisonnalité, ainsi que la sélection
automatique du meilleur modèle par maladie constituent les piliers méthodologiques
du pipeline. La traduction de ces projections en niveaux d'alerte conformes aux
standards OMS/SIMR, assortie de garde-fous quantitatifs et de justifications
transparentes transmises aux autorités, confère à la plateforme \textsc{SAFE CONGO}
une valeur opérationnelle concrète au service de la santé publique congolaise.
